\chapter{Discussion}

\section{Comparison of Spectrometers}
When trying to evaluate the performance of the homebuild prism spectrometer we use OceanOptics flame\autocite{OcOptFlameManual2016} as comparison. 
One should be aware of the core difference between the two spectrometers, being the dispersion method. 
OceanOptics flame used deffractive dispersion with a grating\autocite{PaschottaDiffractionGratings} that provides a near linear angular dispersion. 
Our prism setup relies on the non-linear Sellmeier equation~\eqref{eq:Sellmeier}.

As seen in Figure~\ref{fig:SpectrometerComparison}, the resolution of our prism spectrometer is notably lower. 
Problems in the homemade spectrometer could be spherical abberation\autocite{PaschottaSphericalAberrations}, the use of an achromatic doublet lens could have led to better collimation. 
The use of a narrower slit could have resulted in increased spectral resolution\autocite{PaschottaSpectrometers} and gotten better seperation of peaks, but it also lowers the intensity, so a more intense lamp should have followed this improvement. 

The lambda solver was very unstabil\dots

\section{Solar Temperature}
Our calculated solar temperature of $5564.1 \text{ K} \pm 12.6 \text{ K}$ is remarkably close to the accepted effective temperature at $5770 \text{ K}$ \autocite[p. 574]{LissauerDePater2013}. 
$$\frac{|5564.1 \text{ K} - 5770 \text{ K}|}{5770 \text{ K}} \cdot 100\% \approx 3.6\%$$
A discrepancy of $3.6\%$ is probably fair for a measurement taken with a handheld fiber-optic cable and a portable spectrometer.